% Section 2.4, from lectures/L02/README.md: the goals, the evaluation questions, and the next
% lecture. The README's Reference list, pointing at Appendices C and D, is in the chapter opener.

\needspace{10\baselineskip}
\section{Review}
\label{c2:sec:review}

\subsection*{What you should be able to do now}
After this chapter, you should:
\begin{itemize}
  \item Understand why embedded systems are programmed close to the hardware, in C, without an OS
    underneath.
  \item Be able to read and write bitwise operations (\code{|}, \code{&}, \code{~}, \code{<<},
    \code{>>}) to set, clear, and read individual bits in a register.
  \item Be able to configure a pin as input or output and drive/read it correctly.
\end{itemize}

\needspace{6\baselineskip}
\subsection*{Questions to test yourself}
You should be able to explain:
\begin{enumerate}
  \item What's the difference between \code{DDRx}, \code{PORTx}, and \code{PINx}, and when do you
    write to one versus read from another?
  \item How do you set bit 3 of a register without disturbing the other bits? How do you clear it?
  \item Why does an input pin need its pull-up resistor enabled, and what would you expect to happen
    if you left it disabled with nothing else connected to the pin?
\end{enumerate}

\needspace{6\baselineskip}
\subsection*{Looking ahead}
Chapter~\ref{c3:ch} is about interrupts:
\begin{itemize}
  \item Interrupts: reacting to events instead of polling for them.
  \item External interrupts (INT0/INT1) and Pin Change Interrupts.
\end{itemize}
